\documentclass[a4paper,8pt]{article}

\usepackage[top=0.25in, bottom=0.25in, left=0.4in, right=0.4in]{geometry}
\usepackage{fontspec}
\usepackage{enumitem}
\usepackage{titlesec}
\usepackage[hidelinks]{hyperref}

\setmainfont{Helvetica}
\setlength{\parindent}{0pt}
\setlength{\parskip}{0pt}

\titleformat{\section}{\normalsize\bfseries}{}{0em}{}[\titlerule]
\titlespacing*{\section}{0pt}{3pt}{1pt}

\setlist[itemize]{nosep, topsep=1pt, leftmargin=1.5em, label=\textbullet, itemsep=0pt, parsep=0pt}

\pagestyle{empty}

\begin{document}

% ===== KİŞİSEL BİLGİLER =====
\begin{center}
{\Large\bfseries Emre Pirinç}\\[2pt]
{\small AI Researcher | SAP B1 Consultant | Enterprise AI Solutions | SAP Joule AI | MIS Graduate}\\[2pt]
{\small İstanbul, Türkiye | info@1takimstartuplar.com | 05345905889}\\[1pt]
{\small LinkedIn: linkedin.com/in/emre-pirinc}
\end{center}

\vspace{-2pt}

% ===== ÖN YAZI =====
\section{ÖN YAZI}
{\small Sakarya Üniversitesi Yönetim Bilişim Sistemleri mezunuyum. SAP Business One Uygulama Danışmanlığı alanında stajyerlik sürecimin ardından aktif projelerde görev alarak toplam 1 yıllık uçtan uca SAP B1 danışmanlığı deneyimi kazandım. Staj sürecimin hemen ardından, Temmuz – Şubat döneminde sahada (yerinde) yürütülen projede görev alarak analiz, tasarım, geliştirme ve canlıya geçiş aşamalarında aktif sorumluluk üstlendim. İhtiyaç analizi, proje yönetimi, raporlama ve sistem entegrasyonu süreçlerinde çalışarak uçtan uca proje deneyimi kazandım.

AIFTeam'de görev yaptığım süreçte, müşteri ihtiyaçlarına yönelik çözümler geliştirerek ekip çalışmasıyla projelerin başarıyla tamamlanmasına katkı sağladım. Crystal Reports ile rapor tasarımı ve SAP B1 sistem entegrasyonları üzerine çalışarak şirketlerin operasyonel verimliliğini artırmaya odaklandım. Bu süreç, analitik düşünme, problem çözme ve iş birliği yetkinliklerimi güçlendirdi.

Kariyerime Yapay Zeka Mühendisliği ve AI Danışmanlığı alanında devam etmek istiyorum. Bu doğrultuda AI projelerinde aktif rol almak, üretmek ve değer yaratmak hedefindeyim. Uzun vadeli hedefim; SAP Joule AI çözümleri ve AI destekli kurumsal yapay zeka sistemleri üzerine uzmanlaşarak kendi girişimimi kurmak.

Teknoloji aşığı biriyim. Bir işin doğru yapılması benim için çok önemlidir, bu konuda inatçı olabilirim. Sevdiğim bir işte inanılmaz bir efor ile saatlerce tam odakla çalışabilirim. Maaştan çok iş arkadaşlarım ve gördüğüm muamele benim için önemlidir. Öğrenmediğim ve vaktimi kaybettiğim prosedürler canımı sıkar, sırf yapmak için iş yapmam. İlk geldiğim ortamda gözlemci ve sessiz olurum ama kısa sürede herkesle konuşan, çok soru soran birine dönüşürüm. İnsanlara yardım etmeyi severim. İş dışında FPS oyunlarını, doğa yürüyüşlerini ve bisiklet kullanmayı severim.}

% ===== YETENEKLER =====
\section{YETENEKLER}
{\small SAP Business One, SAP ERP, Crystal Reports, SQL, Python, Microsoft Excel, Project Management, Business Analysis, Artificial Intelligence, Machine Learning, LLM, Data Science, Data Analysis, Scrum, Figma, HTML, CSS, Blockchain, Test Automation, Prompt Engineering, Veritabanı Yönetimi, ISO 9001, Kanban}

% ===== İŞ DENEYİMİ =====
\section{İŞ DENEYİMİ}

\textbf{1 Takım Start Uplar} \hfill {\small Şubat 2026 -- Devam Ediyor}\\
{\small \textit{Kendi İşim} -- İstanbul, Türkiye}
\begin{itemize}
\item {\small Kendi arkadaş çevrem ile kurmuş olduğum ekibimde projeler yapmakta, startup yarışmalarına katılmakta, haftalık ve aylık toplantılar ile proje toplantıları yapmaktayız. 50'ye yakın proje fikri ve aktif olarak 3 proje üzerinde çalışmaktayız.}
\end{itemize}

\vspace{0pt}

\textbf{AIFTeam | SAP Business One Certified Partner} \hfill {\small 1 yıl 1 ay}\\
{\small \textit{İstanbul, Türkiye}}

\vspace{0pt}
{\small \textbf{HERO Ürün Danışmanı} · Tam Zamanlı \hfill Ocak 2026 -- Şubat 2026}\\
{\small Uzaktan}
\begin{itemize}
\item {\small AIFTeam bünyesinde yürütülen Hero projesinde; HMR ve PM modülleri ile mobil uygulamaları içeren entegre portalın uçtan uca süreç yönetimi ve operasyonel verimlilik çalışmalarında aktif rol almaktayım.}
\end{itemize}

\vspace{0pt}
{\small \textbf{SAP B1 Consultant} · Tam Zamanlı \hfill Temmuz 2025 -- Şubat 2026}\\
{\small Ofisten}
\begin{itemize}
\item {\small SAP Business One danışmanlık süreçlerinde analiz, tasarım, geliştirme ve canlıya geçiş aşamalarının uçtan uca yönetiminde görev almaktayım.}
\item {\small Proje Yönetimi \& Analiz: Müşteri ihtiyaç analizlerinin yapılması, Kavramsal Tasarım Dokümanlarının (Blueprint) hazırlanması ve proje süreçlerinin takibi.}
\item {\small Raporlama \& Tasarım: Crystal Reports ile şirket ihtiyaçlarına yönelik özel raporların tasarlanması, form dizaynlarının yapılması ve SQL sorgularının optimize edilmesi.}
\item {\small Kurulum \& Konfigürasyon: Şirket, Server, Client ve SAP B1 kurulumlarının gerçekleştirilmesi ve sistem parametrelerinin uyarlanması.}
\item {\small Entegrasyon \& Geliştirme: Web projeleri ve 3. parti yazılımlar ile SAP B1 entegrasyonlarının yönetimi; add-on süreçlerine destek verilmesi.}
\item {\small Test \& Destek: Kullanıcı Kabul Testleri (UAT) ve canlı kullanım sonrası destek süreçlerinin yürütülmesi.}
\end{itemize}

\vspace{0pt}
{\small \textbf{SAP B1 Consultant} · Stajyer \hfill Şubat 2025 -- Haziran 2025}\\
{\small Ofisten}

\vspace{0pt}

\textbf{Intern} \hfill {\small Temmuz 2024 -- Eylül 2024}\\
{\small \textit{Seratu Medya Reklam Yazılım Danışmanlık · Stajyer} -- Konya, Türkiye}
\begin{itemize}
\item {\small Ortaokul seviyesine yönelik Bilişim ve Teknoloji Tasarım kitaplarının scratch robotik kodlama ve algoritma içeriklerini hazırladım. Görsel tasarım süreçlerinde Stable Diffusion kullanarak yapay zeka destekli içerikler ürettim ve QR kod entegrasyonu ile interaktif ders materyalleri tasarladım.}
\end{itemize}

\vspace{0pt}

\textbf{Paketleme} \hfill {\small Ağustos 2023 -- Ekim 2023}\\
{\small \textit{TTT-auto · Dönemsel} -- Konya, Türkiye}
\begin{itemize}
\item {\small Global ölçekte üretim yapan firmanın fren kaliperi ve tamir takımları paketleme departmanında görev aldım. İhracat standartlarına uygun paketleme süreçlerini yürütürken, akıllı depolama sistemleri (WMS) ve otomasyon teknolojilerini yerinde gözlemleme fırsatı buldum.}
\end{itemize}

\vspace{0pt}

\textbf{TİİBÖT Sosyal Medya Editörü} \hfill {\small Mart 2022 -- Haziran 2022}\\
{\small \textit{TİİBÖT} -- Türkiye}
\begin{itemize}
\item {\small Öğrenci topluluğunun dijital iletişim süreçlerinde gönüllü olarak görev aldım. Etkinlikler için sosyal medya içerikleri hazırladım ve organizasyonlara uzaktan katılarak yönetim ekibine operasyonel destek verdim.}
\end{itemize}

% ===== EĞİTİM =====
\section{EĞİTİM}

\textbf{Sakarya Üniversitesi} \hfill {\small Ekim 2021 -- Haziran 2025}\\
{\small \textit{Lisans, Yönetim Bilişim Sistemleri} -- Not: 3.27}

\vspace{0pt}

\textbf{İstanbul Üniversitesi} \hfill {\small Ekim 2025 -- Temmuz 2027}\\
{\small \textit{Associate's Degree, Computer Programming}}

\vspace{0pt}

\textbf{İzmir Bakırçay Üniversitesi} \hfill {\small Eylül 2019 -- Eylül 2021}\\
{\small \textit{Lisans, Yönetim Bilişim Sistemleri}}

% ===== PROJELER =====
\section{PROJELER}

\textbf{Sakarya Üniversitesi Öğrencilerinin Yapay Zeka Bağımlığı Ve Korkuları} \hfill {\small Kasım 2023 -- Devam Ediyor}\\
{\small TÜBİTAK 2209-A destekli proje. Sakarya Üniversitesi ile ilişkili.}

\vspace{0pt}

\textbf{Anadolu Bakır SAP Business One Canlı Geçiş Projesi} \hfill {\small Temmuz 2025 -- Şubat 2026}\\
{\small AIFTeam | SAP Business One Certified Partner ile ilişkili.}

\vspace{0pt}

\textbf{QDMS Kalite Doküman Yönetimi Sistemi Tasarlanması} \hfill {\small Kasım 2024 -- Aralık 2024}\\
{\small 7 haftalık sprintler boyunca ekip olarak ISO 9001 standartlarına uygun bir kalite yönetim sistemi geliştirdik. Müşteri gereksinimlerini analiz ederek çalışan bir yazılım sunduk. Sakarya Üniversitesi ile ilişkili.}

\vspace{0pt}

\textbf{SMA'YA IŞIK TUT} \hfill {\small Ocak 2024 -- Haziran 2024}\\
{\small SMA ve bağışçıları bir araya getiren ve farkındalık amaçlayan bir mobil uygulamanın teorik çalışması. Gereksinim belirleme, UI/UX tasarımı, proje planlama ve risk yönetimi. Sakarya Üniversitesi ile ilişkili.}

\vspace{0pt}

\textbf{Uber Gereksinim Analizi Yapılması} \hfill {\small Ocak 2024 -- Haziran 2024}\\
{\small Uber iş modelini detaylı analiz ederek gereksinimlerini tanımladık. Use case diyagramları, wireframe, storyboard ve story mapping tekniklerinin kullanımı. Sakarya Üniversitesi ile ilişkili.}

\vspace{0pt}

\textbf{E-Park: Akıllı Otopark Çözümleri YBS Projesi} \hfill {\small Ocak 2021 -- Haziran 2021}\\
{\small Yoğun trafik ve park sorunlarına çözüm getiren, mobil tabanlı, entegre akıllı otopark sistemi projesi. Plaka Tanıma Sistemi (PTS), Masterpass ödeme altyapısı ve Blockchain destekli veri güvenliği. Sakarya Üniversitesi ile ilişkili.}

% ===== LİSANSLAR VE SERTİFİKALAR =====
\section{LİSANSLAR VE SERTİFİKALAR}
{\small
\textbf{SAP Certified -- SAP Business One} -- SAP (Ocak 2026)\\
\textbf{Artificial Intelligence -- 2602 Early Release Series} -- SAP (Ocak 2026)\\
\textbf{Introduction to Large Language Models} -- Google Cloud Training Online (Ocak 2026)\\
\textbf{Insider One AI Weekend Certification} -- Insider One (Şubat 2026)\\
\textbf{Google Cloud Introduction to Data Engineering} -- Istanbul Data Science Academy (Mart 2026)\\
\textbf{CSCON'26 (IEEE Türkiye Computer Society Congress)} -- IEEE Computer Society Türkiye SAC (Mart 2026)\\
\textbf{Veri Bilimi ve Yapay Zeka Zirvesi Datacamp'2025} -- Boğaziçi Üniversitesi (Aralık 2025)\\
\textbf{Python ile Programlamaya Giriş DATACAMP'25} -- Boğaziçi Üniversitesi (Aralık 2025)\\
\textbf{Querying MS SQL} -- Miuul (Şubat 2025)\\
\textbf{Uygulamalarla SQL Öğreniyorum} -- BTK Akademi (Ocak 2025)\\
\textbf{SAP Business One Katılım Sertifikası} -- UpexTech (Ocak 2025)\\
\textbf{Data Analysis in Excel} -- DataCamp (Kasım 2024)\\
\textbf{Excel'e Giriş} -- DataCamp (Kasım 2024)\\
\textbf{Time Series Analysis in Python} -- DataCamp (Kasım 2024)\\
\textbf{Veri Bilimi Zirvesi Datacamp24} -- Boğaziçi Üniversitesi (Kasım 2024)\\
\textbf{Business Analysis and PO with AI Assistance} -- BA-WORKS İş Analizi Hizmetleri (Ağustos 2024)\\
\textbf{Veri Ön İşleme} -- T.C. Cumhurbaşkanlığı Dijital Dönüşüm Ofisi (Mart 2024)\\
\textbf{Zorlu Süreçlerden Geçerken Stresi Yönetmek} -- Sakarya Üniversitesi (Mart 2024)\\
\textbf{Haliç Üniversitesi Bilişim ve İnovasyon Zirvesi} -- Haliç Üniversitesi YBS Kulübü (Nisan 2023)\\
\textbf{Yeni Başlayanlar için BlockChain} -- SAÜ Blockchain Öğrenci Topluluğu (Aralık 2022)\\
\textbf{Boğaziçi DataCamp'22 Veri Bilimi Zirvesi} -- Boğaziçi Üniversitesi (Kasım 2022)\\
\textbf{Blockchain Okuryazarlığı Farkındalık Eğitimi} -- Habitat Derneği (Mart 2022)\\
\textbf{How to Start Girişimcilik Zirvesi} -- Boğaziçi Üniversitesi İşletme ve Ekonomi Kulübü (Mart 2022)\\
\textbf{Boğaziçi DataCamp'21 Veri Bilimi Zirvesi} -- Boğaziçi Üniversitesi (Aralık 2021)
}

% ===== ONUR VE ÖDÜLLER =====
\section{ONUR VE ÖDÜLLER}
{\small
\textbf{Zula Bitexen Türkiye Kupası} -- TESFED \& Bitexen (Ağustos 2021) -- BALEspor Akademi takımı ile Türkiye 2.'liği. IGL olarak görev aldım.\\
\textbf{Zula Süper Lig 6. Sezon} -- Zula Super Lig (Ağustos 2020) -- Çamlıca Espor takımı kadrosunda profesyonel ligde mücadele ettim.\\
\textbf{Paycell \#EvdeKal Zula Turnuvası \#3} -- Paycell (Mayıs 2020) -- Godzilla e-Sports ile şampiyonluk. Turnuva MVP'si -- Nvidia GTX 1650 kazandım.\\
\textbf{Paycell \#EvdeKal Zula Turnuvası \#4} -- Paycell (Mayıs 2020) -- Godzilla e-Sports ile şampiyonluk. IGL olarak görev aldım.\\
\textbf{Paycell \#EvdeKal Zula TDM Turnuvası \#1} -- Paycell / Zula Oyun (Nisan 2020) -- Godzilla e-Sports ile şampiyonluk. Turnuva MVP'si -- Nvidia GTX 1650 kazandım.\\
\textbf{Paycell \#EvdeKal Zula TDM Turnuvası \#2} -- Paycell (Nisan 2020) -- Godzilla e-Sports ile şampiyonluk. IGL olarak görev aldım.\\
\textbf{ininal Online Zula Turnuvası} -- ininal (Nisan 2020) -- Godzilla e-Sports ile şampiyonluk. IGL olarak görev aldım.
}

% ===== DİLLER =====
\section{DİLLER}
{\small Türkçe (Ana Dil), English (A2)}

\end{document}
